\documentclass[Chp3.tex]{subfiles}

\begin{document}


\section{Proofs}
\label{app:DF_R}

\subsection{Proposition 1}

\begin{proof}

\par Let
\[
W^{A}(\{p_s\};U) \equiv \min_{\lambda \in C} \int_S \big(\mathbb{E}_{p_s}[U]\big)\,d\lambda,
\qquad C \subseteq \Delta(S),\ C\neq\emptyset,
\]
and fix any $U$ with $U' > 0$ and $U'' \leq 0$.

\par Suppose that, for every state $s \in S$,
\[
\int_0^z \big[P_s(y)-P_s^*(y)\big]dy \le 0
\quad\text{for all } z\in[0,\bar y],
\]
where
\[
P_s(y)=\int_0^y p_s(t)\,dt,
\qquad
P_s^*(y)=\int_0^y p_s^*(t)\,dt.
\]

\par For each fixed state $s$, this is exactly the single-state second-order stochastic dominance condition. Hence,
\[
\mathbb{E}_{p_s}[U] \ge \mathbb{E}_{p_s^*}[U]
\quad\text{for every } s\in S.
\]

\par Let $\lambda\in C$ be arbitrary. Since $\lambda$ is a probability measure on $S$,
\[
\int_S \mathbb{E}_{p_s}[U] \, d\lambda
\ge
\int_S \mathbb{E}_{p_s^*}[U] \, d\lambda.
\]
Because this inequality holds for every $\lambda\in C$, taking minima over $C$ gives
\[
\min_{\lambda\in C}\int_S \mathbb{E}_{p_s}[U] \, d\lambda
\ge
\min_{\lambda\in C}\int_S \mathbb{E}_{p_s^*}[U] \, d\lambda.
\]
Therefore,
\[
W^{A}(\{p_s\};U) \ge W^{A}(\{p_s^*\};U).
\]
Since $U$ was arbitrary in the class $\{U: U'>0,\ U''\le 0\}$, the claim follows.

\end{proof}

\begin{remak}
The argument above establishes a sufficient condition. In general, the converse fails under a maxmin aggregator: the minimizing prior may assign zero weight to states where state-by-state dominance fails. An equivalence requires stronger assumptions (for example, comparisons for every $\lambda$ in a sufficiently rich set of priors that identifies each state).
\end{remak}

\subsection{Single-state benchmark}

\begin{prop}
For a single-state problem with distributions $f$ and $f^*$ on $[0,\bar y]$, define
\[
F(y)=\int_0^y f(t)\,dt,
\qquad
F^*(y)=\int_0^y f^*(t)\,dt.
\]
Then
\[
\mathbb{E}_{f}[U] \ge \mathbb{E}_{f^*}[U]
\quad\text{for all } U\text{ with }U'>0,\ U''\le 0
\]
if and only if
\[
\int_0^z \big[F(y)-F^*(y)\big]dy \le 0
\quad\text{for all } z\in[0,\bar y].
\]
Moreover, strict preference for all such $U$ requires strict inequality for some $z\in(0,\bar y)$.
\end{prop}

\subsection{Section 3.6 The Analogy between Income and Wealth Inequality Aversion}
\label{app:objective-benchmark}

\par To understand the problem, I begin with describing the model which proposes measures of income inequality using the choice under objective uncertainty literature. 

\subsubsection*{Analytical framework}

\par Denote the set of outcomes $Y = [0,\bar{y}]$. Then, the choice set is given by:

$$ L = \bigg\{ p: 2^{[0,\bar{y}]} \to \mathbb{R} \hspace{2mm} | \hspace{2mm} p(\cdot) \hspace{2mm} \text{is an income frequency distribution}  \bigg\}. $$

\par We may define a binary relation over both sets $Y,L:$

$$ \succsim_{y} \subseteq [0,\bar{y}] \times [0,\bar{y}] \subseteq \mathbb{R} \times \mathbb{R} $$

$$ \succsim \; \subseteq L \times L. $$

\par Notice that $\succsim_y$ may be represented by a real-valued utility function. This is the social welfare $U(y)$ which is a key object of analysis in this paper. 

\subsubsection*{Axiomatization}

\begin{ax.V}[Weak Order]
$\succsim$ is complete and transitive.
\end{ax.V}

\begin{ax.V}[Continuity]
For every $p(\cdot)$, $p^*(\cdot)$, $p'(\cdot)$ $\in L$, if $p(\cdot) \succ p^*(\cdot) \succ p'(\cdot) $, there exists $\alpha, \beta \in (0,1)$ such that 
$$ \alpha p + (1-\alpha)p' \succ p^* \succ \beta p + (1-\beta) p' .$$ 
\end{ax.V}

\begin{ax.V}[Independence]
For every $p,p^*,p' \in L$ and every $\alpha \in (0,1)$ such that 
$$ p \succsim p^* \iff \alpha p + (1-\alpha)p' \succsim \alpha p^* + (1-\alpha)p'. $$
\end{ax.V}

\subsubsection*{Expected utility representation of the ranking over income distributions}

\par The ranking over income distributions will be represented using the vNM-expected utility representation. That is,
$$ p(y) \sim \int_{0}^{\bar{y}} U(y) p(y) dy \equiv W^{O}. $$

\par Formally, note the slight modification of the vNM-EU theorem in this setting of ranking income distributions:

\begin{tm}
$\succsim \subseteq L \times L$ satisfies $(V1)$ \textbf{weak order}, $(V2)$ \textbf{continuity}, $(V3)$ \textbf{independence} if and only if there exists $U: Y \to \mathbb{R}$ such that, for every $p(y), p^*(y) \in L,$
$$ p(y) \succsim p^*(y) \iff \int_{0}^{\bar{y}} U(y) p(y) dy \geq \int_{0}^{\bar{y}} U(y) p^*(y) dy. $$
\end{tm}

\subsubsection*{3.6.1 Ratio measure $D$}

Recall the total social welfare under the two choice models: choice under objective uncertainty and choice under ambiguity. 
\[
W^{O} \equiv \int_0^{\bar y} U(y)\,p(y)\,dy
\]
\[
W^{A} \equiv \min_{\lambda\in C}\int_S \mathbb{E}_{p_s}[U] \, d\lambda.
\]
Then the ratio measures are
\[
D^{O} = \frac{W^{O}}{U(\mu)},
\qquad
D^{A} = \frac{W^{A}}{U(\mu)}.
\]


\subsubsection*{3.6.2 Equally distributed equivalent measure}

Define the EDE wealth levels by
\[
U\big(w_{EDE}^O\big)=W^{O} = \int_0^{\bar y} U(y)\,p(y)\,dy,
\qquad
U\big(w_{EDE}^A\big)= W^{A} = \min_{\lambda\in C}\int_S \mathbb{E}_{p_s}[U] \, d\lambda.
\]
Hence,
\[
w_{EDE}^O = U^{-1}(W^{O}),
\qquad
w_{EDE}^A = U^{-1}(W^{A}).
\]
The corresponding inequality measures are
\[
I^{O} = 1-\frac{w_{EDE}^O}{\mu},
\qquad
I^{A} = 1-\frac{w_{EDE}^A}{\mu}.
\]

\subsubsection*{3.6.3 Closed form under homothetic preferences}

For
\[
U(y)=
\begin{cases}
\alpha + \beta\dfrac{y^{1-\epsilon}}{1-\epsilon}, & \epsilon\neq 1,\\
\alpha + \beta\ln y, & \epsilon=1,
\end{cases}
\qquad \beta>0,
\]
we obtain the objective and ambiguity formulas by substituting the corresponding welfare aggregator.

If $\epsilon\neq 1$:
\[
w_{EDE}^O=\left[\frac{1-\epsilon}{\beta}\big(W^{O}-\alpha\big)\right]^{\frac{1}{1-\epsilon}},
\qquad
w_{EDE}^A=\left[\frac{1-\epsilon}{\beta}\big(W^{A}-\alpha\big)\right]^{\frac{1}{1-\epsilon}},
\]
\[
I^{O} = 1-\frac{1}{\mu}\left[\frac{1-\epsilon}{\beta}\big(W^{O}-\alpha\big)\right]^{\frac{1}{1-\epsilon}},
\qquad
I^{A} = 1-\frac{1}{\mu}\left[\frac{1-\epsilon}{\beta}\big(W^{A}-\alpha\big)\right]^{\frac{1}{1-\epsilon}}.
\]

If $\epsilon=1$:
\[
w_{EDE}^O=\exp\!\left(\frac{W^{O}-\alpha}{\beta}\right),
\qquad
w_{EDE}^A=\exp\!\left(\frac{W^{A}-\alpha}{\beta}\right),
\]
\[
I^{O} = 1-\frac{1}{\mu}\exp\!\left(\frac{W^{O}-\alpha}{\beta}\right),
\qquad
I^{A} = 1-\frac{1}{\mu}\exp\!\left(\frac{W^{A}-\alpha}{\beta}\right).
\]

\par So the extension from objective uncertainty to ambiguity is algebraically the replacement
\[
W^{O} \longrightarrow W^{A},
\]
with the EDE inversion and inequality-index algebra unchanged.


\subsection{Gross-assets levels by race}
\label{app:gross-levels-race}

\par Race-specific \textit{mean} net wealth and mean gross assets by year appear in the main text as Figure~\ref{fig:mean-net-gross-by-race}, immediately before the Black--White median-gap diagnostic (Figure~\ref{fig:gap-net-gross}).

\subsection{Restricted-sample median-gap diagnostic}
\label{app:restricted-gap-diagnostic}

\begin{figure}[H]
\centering
\includegraphics[width=0.92\linewidth]{\subfix{../Figures/03_gap_diagnostic_positive_only.png}}
\caption{Black-White median wealth differential in net wealth and gross assets, positive-only samples (SCF 1989--2022).}
\label{fig:gap-net-gross-restricted}
\end{figure}

\par Even after imposing positive-only restrictions, both median-gap series still display visible time variation and an upward tendency over the sample. This provides a useful diagnostic benchmark for interpreting the unusually stable, high restricted-sample model-based inequality indices in the main text.

\subsection{Shift constant and mass at non-positive wealth}
\label{app:shift-mass-diagnostic}

\begin{figure}[H]
\centering
\includegraphics[width=0.95\linewidth]{\subfix{../Figures/03_shift_constant_and_nonpos_share.png}}
\caption{Left: year-specific shift $c_{t,m}$ applied in the shifted full-sample treatment. Right: weighted fraction of households with non-positive raw wealth for each measure (before shifting).}
\label{fig:shift-mass-diagnostic}
\end{figure}

\par Together with the median-gap figures, these panels are meant as a simple check on how the sample is affected when households with non-positive wealth are kept in the shifted full-sample treatment (by adding a year-specific constant so utility is well-defined) rather than dropped as in the positive-only restriction. The left panel shows how large that shift had to be each year; the biggest values line up with 2010--2016, which is the period when balance-sheet stress and indebtedness were most salient in the aggregate SCF. The same window shows up in Figure~\ref{fig:mean-net-gross-by-race} as an episode where mean gross assets are comparatively soft in levels while mean net wealth is also compressed---the configuration you expect when liabilities are weighing heavily on net positions. The right panel shows that the weighted share of observations at or below zero in raw wealth is largest in those years as well, for both net wealth and gross assets, so the mechanical ``mass at zero or below'' lines up with when the shift has to do the most work.

\subsection{Computational benchmark: the linear case (epsilon = 0)}
\label{app:eps0-benchmark}

\par When $\epsilon=0$, the CRRA class collapses to linear utility, $U(y)=\alpha+\beta y$ with $\beta>0$. Under this case, the objective index satisfies
\[
I^{O}=0
\]
by construction, since the equally distributed equivalent equals the mean. This is the expected no-inequality-aversion benchmark.

\par In contrast, the ambiguity index remains sensitive to cross-state asymmetry because maxmin aggregation selects the lower state-level expected utility. In the SCF computations, this yields positive $I^{A}$ values even at $\epsilon=0$, with notably larger values in the restricted positive-only samples than in most shifted full-sample years.

\par The $\epsilon=0$ output is therefore used as a calibration benchmark for the main text: it isolates the linear, no-inequality-aversion case while preserving the ambiguity-vs-objective comparison.


\end{document}
